\BourbakiExercises{习题}

\begin{enumerate}
\def\labelenumi{\arabic{enumi})}
\tightlist
\item
  设 \(D\) 是一个体。
\end{enumerate}

\begin{enumerate}
\def\labelenumi{\alph{enumi})}
\item
  设 E 是 D 的一个真子域，其乘法群 \(E^{*}\) 在 \(D^{*}\) 中有有限指标。证明 D 是有限的（对 E 的互异元素的一个序列 \((a_{n})\) 与 D---E 的一个元素 x，考虑在模 \(E^{*}\) 意义下元素 \(x + a_{n}\) 的类）。。
\item
  证明 D 中只有有限多个共轭元的元素属于 D 的中心 Z（对该元素在 D 中的交换子应用 a)，它是一个体）。
\item
  由此推出，Z{[}X{]} 中在 D-Z 中有一个根的多项式在 D-Z 中有无穷多个根。
\end{enumerate}

¶ 2) 设 A 是一个 K-代数，m 是一个整数。假设 A 的每个元素在 K 上是次数 \(\leqslant m\) 的代数元。

\begin{enumerate}
\def\labelenumi{\alph{enumi})}
\item
  证明：若 A 是一个本原环，则存在一个整数 \(r \leqslant m\) 与一个体 D（其元素在 K 上的次数都 \(\leqslant m\)），使得 A 同构于矩阵代数 \(\mathbf{M}_{r}(\mathrm{D})\)。（若 A 不是单的，则由稠密性定理推出，对每个整数 n，存在 A 的一个子代数 \(A_{n}\) 与一个满射同态 \(A_{n} \to M_{n}(\mathrm{D})\)。注意代数 \(\mathbf{M}_{n}(\mathrm{D})\) 包含满足 \(X^{n-1} \neq 0\) 与 \(X^{n} = 0\) 的元素 X。）
\item
  设域 K 是无限的。证明代数 \(A/\Re(A)\) 是半单的且至多有 m 个单分量（注意代数 \(K^{n}\) 包含 K 上次数为 n 的元素；由此推出 A 至多有 m 个极大双边理想，并用 a) 结束证明）。证明：若 K 还是完满的，则 \(A/\Re(A)\) 在 K 上有有限秩。
\end{enumerate}

¶ 3) 设 K 是一个交换域。我们说 K-代数 A 是代数的，如果它的所有元素在 K 上都是代数的。每个有限秩代数都是代数的；代数代数的每个子代数与商代数都是代数的。代数 K-代数的根是一个诣零理想（VIII, 第 166 页，习题 5）。

\begin{enumerate}
\def\labelenumi{\alph{enumi})}
\item
  证明代数代数是一个伪正则环（VIII, 第 178 页，习题 4；注意对每个 \(x \in \mathrm{A}\)，存在一个元素 \(y \in \mathrm{K}[x]\) 与一个整数 \(k\)，使得 \(x^k = x^{k+1}y\)）。
\item
  设 A 是一个除 0 外没有其他幂零元素的代数 K-代数。证明 A 同构于诸体的一个乘积的子代数。（注意 A 的每个商代数具有同样的性质。像习题 2, a) 那样推理，证明若 A 是本原的，则它是一个体。在一般情形，使用 VIII, 第 168 页，习题 12。）
\end{enumerate}

¶ 4) 设 K 是一个有限域，A 是一个代数 K-代数（习题 3）。

\begin{enumerate}
\def\labelenumi{\alph{enumi})}
\item
  设 A 是一个体。证明 A 是交换的。（设 \(x\) 是 A 的一个非中心元素。考虑体 \(Z(x)\)（其中 \(Z\) 是 A 的中心），并证明存在一个元素 \(y \in A\) 使得 \(yxy^{-1} = x^r\)（且 \(x^r \neq x\)）。通过考虑 A 的子域 \(K(x,y)\) 推出一个矛盾。）
\item
  设 A 不含任何非零的幂零元素。证明代数 A 是交换的（用 a) 与习题 3）。
\end{enumerate}

\begin{enumerate}
\def\labelenumi{\arabic{enumi})}
\setcounter{enumi}{4}
\item
  设 K 是一个有限域，m 是一个整数，A 是一个无根的 K-代数，其元素的次数都 \(\leqslant m\)。证明存在 K 的一个有限扩张 L，使得 A 同构于一个乘积 \((\mathbf{M}_{r}(\mathrm{L}))^{\mathrm{I}}\) 的子代数，且该乘积的投影都是单代数（参看习题 2, a)）。也证明其逆。
\item
  设 A 是一个环，使得对每个 \(a \in \mathrm{A}\)，存在一个整数 \(n(a) > 1\) 满足 \(a^{n(a)} = a\)。证明 A 同构于有限多个有限域上的既约交换代数代数的一个乘积。（注意 A 被某个整数 \(n\) 零化，然后注意对每个素数 \(p\)，加群 A 的 \(p\) -分量 \(\mathrm{A}_p\) 是一个被 \(p\) 零化的双边理想。应用习题 4, b)。）也证明其逆。
\item
  \begin{enumerate}
  \def\labelenumii{\alph{enumii})}
  \tightlist
  \item
    证明：若域 K 具有性质 \((\mathrm{C}_{1})\)（VIII, 第 357 页），则 K 的每个代数扩张也具有该性质。（化归为 K 的一个次数为 r 的有限扩张 L 的情形。证明：若 F 是 L-向量空间 V 上一个次数为 d 的齐次多项式函数，则函数 \(N_{L/K} \circ F\) 在 K-向量空间 V 上是次数为 rd 的齐次多项式。）
  \end{enumerate}
\end{enumerate}

\(^{*}b)\) 从现在起设 K 是代数闭的。设 \(F_{1},\ldots,F_{n}\) 是 m 个变量的、次数 \textgreater0 的齐次多项式。证明：若 n \textless{} m，则存在一个 \(a \neq 0\)（属于 \(K^{n}\)），使得 \(F_{1}(a) = \cdots = F_{n}(a) = 0\)（用 AC, VIII, §2, n° 4, 第 19 页，定理 3 的推论 1 与 AC, VIII, §3, n° 1, 第 24 页，命题 2）。c) 由 b) 推出，域 \(K(X)\) 具有性质 \((C_{1})\)。

（设 f 是 n 个变量的、次数为 d 的齐次多项式，其系数是 X 的多项式，设 \(\nu\) 是这些系数的次数的最大值。设 N 是一个整数，\(G_{1},\ldots,G_{n}\) 是 X 的、次数 \(\leqslant N\) 的多项式。证明等式 \(\mathrm{F}(\mathrm{G}_{1}(\mathrm{X}),\ldots,\mathrm{G}_{n}(\mathrm{X}))=0\) 等价于 \(Nd+\nu+1\) 个次数 \textgreater0 的齐次多项式在 \(n(N+1)\) 个系数（即 \(G_{i}\) 的系数）上全为零。）

\(d\) ) 由 \(a\) ) 与 \(c\) ) 推出，一个代数闭域的、超越次数 \(\leqslant 1\) 的扩张具有性质 \((\mathrm{C}_1)\)（``曾定理''）。*

